\chapter{The Workflow Model}
\label{ch:workflow-model}

\newthought{Before the backends, the shape.} A task card can declare many calculators and
operas; \Jtwo\ turns them into a single \emph{execution plan} --- a layered sequence that
every Sample follows inside its Worker. This chapter explains how layers are derived from
dependencies, what runs concurrently, and how the plan travels.

\section{From modules to layers}
\label{sec:layers}

Each calculator module may declare the modules it depends on:

\begin{yamlwin}
Calculators:
  Modules:
    - name: SPheno            # no dependencies -> layer 0
    - name: HiggsTools
      required_modules: [SPheno]     # -> layer 1
    - name: MicrOmegas
      required_modules: [SPheno]     # -> layer 1 (same layer as HiggsTools)
    - name: Combine
      required_modules: [HiggsTools, MicrOmegas]   # -> layer 2
\end{yamlwin}

Layer assignment is an iterative topological sort over \yamlkey{required\_modules}: a
module's layer is one past the deepest of its dependencies, and modules with no mutual
dependency land in the same layer. \textbf{Same layer = eligible to run concurrently}
inside one Sample.

Two lenient rules keep the sort from being a foot-gun: a dependency name that matches no
module is ignored rather than an error, and a dependency \emph{cycle} falls back to
running the entangled remainder together in one final layer instead of refusing the card.

\begin{jwarn}
Both fallbacks are silent. If your layer structure looks wrong --- things running
together that should be ordered --- check \yamlkey{required\_modules} spellings first: a
typo there does not error, it just flattens your ordering.
\end{jwarn}

\section{The full plan: calculators, then operas, then likelihood}
\label{sec:plan-order}

The complete execution plan for every Sample is assembled deterministically:

\begin{enumerate}
  \item \textbf{Calculator layers} as derived above, starting at layer 0.
  \item \textbf{Operas}, offset to begin after the last calculator layer, in card
        order. Operas therefore always see every calculator's outputs.
  \item \textbf{One likelihood step}, appended as the terminal layer.
\end{enumerate}

For the example above, with one opera \code{L}:

\begin{quote}
\sffamily\footnotesize
[\,SPheno\,@\,0\,]\ \ [\,HiggsTools\,@\,1,\ MicrOmegas\,@\,1\,]\ \
[\,Combine\,@\,2\,]\ \ [\,L\,@\,3\,]\ \ [\,LogLikelihood\,@\,4\,]
\end{quote}

The plan is data-independent: built once from the card, serialized as plain
dictionaries, and shipped to Redis inside every task. Workers iterate it layer by layer
--- \emph{sequential across layers, concurrent within a layer}.

\section{Concurrency inside a Sample}
\label{sec:sample-concurrency}

Within one layer, each calculator run is submitted to the Worker's subprocess scheduler
and waits for its own \emph{global} concurrency slot
(\yamlkey{Calculators.Pools}, Chapter~\ref{ch:calculators}). Three consequences:

\begin{itemize}
  \item Two independent calculators overlap their external run time --- the main win of
        declaring accurate dependencies.
  \item Global caps still hold: eight Workers each running a same-layer
        \code{MicrOmegas} step still respect \code{MicrOmegas}'s total pool.
  \item The widest layer sizes the Worker's scheduler; nothing else needs tuning.
\end{itemize}

Operas run in-process and sequentially, in card order, after all calculators.

\section{Where observables flow}
\label{sec:observable-flow}

The observables dictionary of Section~\ref{sec:observables} threads through the plan:
each calculator's \yamlkey{output} spec and each opera's \yamlkey{output} list merge into
it as their step completes; every later step's expressions and input specs see everything
already merged. Within one concurrent layer, sibling outputs are merged as each finishes
--- do not rely on sibling ordering; if module B needs module A's output, that is by
definition a dependency, so declare it.

\section{Seeing the plan: the flowchart}
\label{sec:flowchart}

Every run exports a semantic flowchart of the workflow --- nodes for parameters,
calculators, operas, and observables, edges from the actual expression and I/O
dependencies --- under the project's \file{images/<scan>/} directory, and \JPlot\ can
render it directly (Chapter~\ref{ch:jarvisplot}). Likelihood nodes are omitted by
default to keep the graph readable.

\begin{jtip}
The flowchart is the fastest sanity check of a new card: if an edge you expect is
missing, the expression or \yamlkey{entry} that should create it references a name that
does not exist --- typically a typo that would otherwise surface as a failed Sample at
runtime.
\end{jtip}
